\begin{corollary}[3.1.4]
\label{IV.2.3.1.4-fr}
Sean $X$ un preesquema localmente noetheriano, $\mathcal F$ un
$\mathcal O_X$-módulo coherente. Entonces los ciclos primos maximales asociados
a $\mathcal F$ son las componentes irreducibles de
$\operatorname{Supp}(\mathcal F)$, y los puntos genéricos de estas componentes
son los puntos $x\in X$ tales que $\mathcal F_x$ sea un
$\mathcal O_x$-módulo de longitud finita y $\ne0$.
\end{corollary}

En efecto, si $x$ es el punto genérico de una de las componentes irreducibles
$Z$ de $\operatorname{Supp}(\mathcal F)$, se sigue de la equivalencia de
\emph{a)} y \emph{c$'$)} en (3.1.3) que $x$ pertenece a
$\operatorname{Ass}(\mathcal F)$, y $Z$ es un ciclo primo asociado a
$\mathcal F$, necesariamente maximal en virtud de (3.1.1.1)~; la recíproca se
sigue trivialmente de (3.1.1.1)~; por último, como la última afirmación es
evidentemente local, se sigue de Bourbaki, \emph{Alg. comm.}, cap.~IV, \S~2,
núm.~5, cor.~2 de la prop.~7.

\begin{corollary}[3.1.5]
\EGAReaderTarget{ega4.c04.cor.3-1-5}{IV.2.3.1.5-fr}
Sean $X$ un preesquema localmente noetheriano, $\mathcal F$ un
$\mathcal O_X$-módulo cuasicoherente. Para que $\mathcal F=0$, es necesario y
suficiente que $\operatorname{Ass}(\mathcal F)=\varnothing$.
\end{corollary}

Como la cuestión es local, se reduce al caso en que $X$ es afín, y la
conclusión se sigue de inmediato de (3.1.2) y de Bourbaki, \emph{Alg. comm.},
cap.~IV, \S~1, núm.~1, cor.~1 de la prop.~2.

